\documentclass[10pt]{article}
\usepackage[margin=0.85in]{geometry}
\usepackage{amsmath,amssymb,amsthm}
\usepackage[hidelinks]{hyperref}
\usepackage{microtype}

\newtheorem{theorem}{Theorem}
\newtheorem{corollary}{Corollary}
\newcommand{\D}{\mathbb D}
\newcommand{\M}{\mathcal M}
\newcommand{\norm}[1]{\lVert #1\rVert_\infty}

\title{A Literature-Implied Positive Subcase of the\\
Uniform Interpolating-Blaschke Approximation Problem}
\author{Run \texttt{fa\_banach\_001} --- lane 13}
\date{9 August 2026}

\begin{document}
\maketitle

\noindent\textbf{Status.}
\emph{Literature-implied answer (partial subcase).}  The implication below is
an agent-identified combination of known theorems.  It is not presented as a
new full solution.

\section*{Source question}

Chalendar, Gorkin, and Partington ask on PDF page 4 of
\cite{CGP2015}:
\begin{quote}
Can every Blaschke product be approximated (uniformly) by an interpolating
Blaschke product?
\end{quote}
Here ``uniformly'' means in the $H^\infty$ norm on the unit disk.  Xin and Hou
still describe the corresponding density problem for the full inner-function
space as open in their 2023 introduction \cite{XinHou2023}.  The two
formulations are equivalent at the level of density because Frostman's theorem
uniformly approximates every inner function by Blaschke products.

\section*{The literature bridge}

For $a\in\D$, put
\[
 \varphi_a(w)=\frac{a-w}{1-\overline a w}.
\]
Following Mortini and Nicolau, let $\M$ be the class of inner functions
$\Theta$ such that $\varphi_a\circ\Theta$ is a Carleson--Newman Blaschke
product for every $a\ne0$.  Thus every such shift is a finite product of
interpolating Blaschke products.  Marshall and Stray prove that a finite
product of interpolating Blaschke products can be approximated arbitrarily
closely in $H^\infty$ norm by a single interpolating Blaschke product; see the
proof of their Lemma 1, especially journal pages 494--495
\cite{MarshallStray1996}.  The modern paper \cite{BNOT2025} records the class
$\M$ in equation (2), proves that every SIP inner function lies in $\M$
(Corollary 1.6(a)), and supplies useful examples and geometric subclasses.

\begin{theorem}
Every inner function $\Theta\in\M$ is a uniform limit in $H^\infty$ of
interpolating Blaschke products.  More quantitatively, for every $a\in\D$
with $a\ne0$, the function
\[
 F_a=\frac{\Theta-a}{1-\overline a\Theta}
\]
is Carleson--Newman and
\begin{equation}\label{eq:error}
 \norm{F_a-\Theta}\le \frac{2|a|}{1-|a|}.
\end{equation}
\end{theorem}

\begin{proof}
Since $F_a=-\varphi_a\circ\Theta$, the defining property of $\M$ says that
$F_a$ is a Carleson--Newman Blaschke product, up to the harmless unimodular
constant $-1$.  Direct subtraction gives
\[
 F_a-\Theta
 =\frac{-a+\overline a\Theta^2}{1-\overline a\Theta}.
\]
Because $|\Theta(z)|\le1$ on $\D$, the numerator has modulus at most $2|a|$
and the denominator has modulus at least $1-|a|$.  This proves
\eqref{eq:error}.

Fix $\varepsilon>0$.  Choose $a\ne0$ so small that
$2|a|/(1-|a|)<\varepsilon/2$.  Since $F_a$ is a finite product of
interpolating Blaschke products, the Marshall--Stray theorem provides a single
interpolating Blaschke product $I$ with
$\norm{I-F_a}<\varepsilon/2$.  The triangle inequality now yields
\[
 \norm{I-\Theta}
 \le \norm{I-F_a}+\norm{F_a-\Theta}<\varepsilon.
\]
\end{proof}

\begin{corollary}
Every Blaschke product whose zeros lie in a Stolz angle is uniformly
approximable by interpolating Blaschke products.  This includes Blaschke
products which are not Carleson--Newman.
\end{corollary}

\begin{proof}
Borichev, Nicolau, Ouna\"ies, and Thomas record that every Blaschke product
whose zeros lie in a Stolz angle belongs to $\M$ \cite[Section 6.5]{BNOT2025}.
Apply the theorem.  Strictness is witnessed by their explicit example with
zeros $1-2^{-j}$ of multiplicity $j$, $j\ge1$: it belongs to $\M$ but is not
Carleson--Newman \cite[Proposition 6.11 and the paragraph following it]{BNOT2025}.
\end{proof}

\section*{Identification and limitations}

The proof is complete, but its inputs are literature theorems and the link is
a direct identification.  The supporting authors do not explicitly advertise
this exact implication as an answer to the question in \cite{CGP2015}; hence
the provenance is ``literature implied,'' not ``literature already answered.''
The result reaches beyond the obvious Carleson--Newman target class, as the
Stolz-angle example shows, but it does not establish that an arbitrary
Blaschke product lies in $\M$.  The full density question therefore remains
open.

\section*{Search record}

The exact question and close variants were searched through arXiv and primary
publisher sources using the phrases ``every Blaschke product,'' ``uniform
approximation,'' ``interpolating Blaschke product,'' ``Frostman shift,'' and
``Carleson--Newman.''  The search located the 2006 modulus theorem of Hjelle
and Nicolau, the 2023 open-status statement of Xin and Hou, the
Marshall--Stray approximation theorem, and the 2025/2026 SIP--$\M$ results.
No full answer was found.  The theorem above is the direct combination that
survived proof checking.

\begin{thebibliography}{9}

\bibitem{CGP2015}
I. Chalendar, P. Gorkin, and J. R. Partington,
\emph{Inner functions and operator theory}, arXiv:1503.05461 (2015).

\bibitem{MarshallStray1996}
D. E. Marshall and A. Stray,
\emph{Interpolating Blaschke products},
Pacific J. Math. \textbf{173} (1996), 491--499.

\bibitem{BNOT2025}
A. Borichev, A. Nicolau, M. Ouna\"ies, and P. J. Thomas,
\emph{Sharp invertibility in quotient algebras of $H^\infty$},
arXiv:2501.05143 (2025).

\bibitem{XinHou2023}
Y. Xin and B. Hou,
\emph{A note on connectedness of Blaschke products},
arXiv:2302.00830 (2023).

\end{thebibliography}

\end{document}
